% Requires amsmath, amssymb and booktabs; data/refinement_rows.tex is generated evidence.
\section{Results, validation and discussion}
\label{sec:results-validation}

We first examine the square in detail, using Section 4's solver results
and additional square-only refinement and load-sensitivity checks with
the same implementation~\cite{nla-square-code}. We then report the
comparisons across objects, loads and supports designed in Section 5,
computed in the accompanying notebook.

\subsection{Deformation under mesh refinement}
We keep Section 4's material, full left clamp and uniform right-edge
traction fixed. Table~\ref{tab:results-refinement} compares structured
meshes with $n$ cells per side, retaining the same diagonal in every
cell. These are separate from Section 5's proposed Triangle meshes.
Each physical measurement uses a checked sparse-LU solution.

\begin{table}[htbp]
 \centering
 \caption{New square-only refinement results from
 \texttt{generate\_evidence.py}. Here $\bar u_x$ is the right-edge
 horizontal displacement averaged by edge length; $u_x^m,u_y^m$ are
 the components at the fixed midpoint $(1,0.5)$ m. Displacements are
 in mm. $\kappa_2(A)$ is estimated from checked extreme eigenpairs.}
 \label{tab:results-refinement}
 \begin{tabular}{rrrrrr}
\toprule
$n$ & Free unknowns & $\bar u_x$ & $u_x^m$ & $u_y^m$ & $\kappa_2(A)$\\
\midrule
5 & 60 & 0.491352 & 0.489202 & 0.010710 & $6.72\times10^{2}$ \\
10 & 220 & 0.492877 & 0.490810 & 0.003769 & $2.47\times10^{3}$ \\
20 & 840 & 0.493497 & 0.491599 & 0.001255 & $9.10\times10^{3}$ \\
40 & 3280 & 0.493744 & 0.491900 & 0.000418 & $3.45\times10^{4}$ \\
\bottomrule
\end{tabular}

\end{table}

The mean extension increases towards about $0.494$ mm, while its
successive changes shrink from $0.00153$ to $0.000620$ and $0.000247$ mm.
The midpoint's vertical movement also decreases. We integrate the edge
trace rather than average nodes, and interpolate at the fixed midpoint
when it is not a node. Absolute differences avoid dividing by its small
vertical component. Every reported mesh difference exceeds the summed
algebraic uncertainty estimates from the two solutions by more than
$10^6$. Those estimates use the actual observation weights, residuals and
smallest eigenvalues; they are not certified bounds. The trend supports
reduced mesh dependence of these displacements, not a known continuum error.

\subsection{Solver behaviour and conditioning}
Section 4's recorded timings show why iteration count alone is not a
speed ranking. Diagonal PCG used 226 steps against CG's 242, but their
minimum setup-inclusive times were 4.99 and 4.31 ms, respectively.
Full GMRES used 228 steps yet took 146 ms. These measurements follow Section 3's protocol; small timing
differences do not establish a general ranking.

Meanwhile, conditioning worsens under refinement despite the settling
displacements. At $n=20$, symmetric diagonal scaling
$P^{-1/2}AP^{-1/2}$ reduces the estimated condition number from 9104 to
8029, where $P=\operatorname{diag}(A)$ as in Section 3. This modest
improvement does not guarantee faster PCG.
The spectral diagnostics also explain the failures. Ordinary Jacobi's
iteration matrix has spectral radius about 1.96, meaning that some error
components grow. Damping by 0.5 reduces this to about 0.999816, below
one but close enough to predict slow asymptotic decay~\cite{saad-iterative}.
This agrees with the recorded divergence of Jacobi and budget exhaustion
of damped Jacobi. Steepest descent also exhausted the budget. Their
runtimes are not times to the requested accuracy. Among accepted runs,
the largest relative difference from the same-mesh LU reference was
$1.74\times10^{-7}$, for Gauss--Seidel, despite its residual being below
$10^{-8}$. Residual and displacement error are different quantities.

\subsection{Sensitivity to small load changes}
At the fixed 20-cell mesh, we compare two load changes of equal relative
size, $\lVert\delta b\rVert_2/\lVert b\rVert_2=10^{-3}$.
The first scales the original load. The second redistributes horizontal
traction along the right edge, with its increment proportional to
$2y/L-1$, where $L=1$ m is the edge height. Its integrated added force
is zero, but its moment changes. The coefficient is chosen to give the
stated load-vector norm; edge forces are integrated consistently.

We solve both changed systems with the same matrix and measure
\begin{equation}
 G=\frac{\lVert\delta x\rVert_2/\lVert x\rVert_2}
          {\lVert\delta b\rVert_2/\lVert b\rVert_2},
 \label{eq:results-amplification}
\end{equation}
where $x$ is the baseline free-displacement vector and $\delta x$ its
change. Uniform scaling gives $G\approx1$, as linearity predicts, while
redistribution gives $G\approx1.77$. Both are far below
$\kappa_2(A)\approx9104$: conditioning describes possible amplification,
not what every load direction produces. Neither test searches for the worst case.
The Euclidean norm here compares free nodal components on one fixed mesh,
not displacement fields on different meshes.

\subsection{Other objects, loads and supports}
Table~\ref{tab:results-objects} compares the three objects of Section 5
at the middle mesh-area limit ($0.00125$ m$^2$). All 36 cases passed
every acceptance check of Section 5.

\begin{table}[htbp]
 \centering
 \caption{Objects under the baseline clamp and traction, middle mesh.
 $\bar u_x$ is the mean right-edge horizontal displacement and $u_x^m$
 the horizontal displacement at the right-edge midpoint, both in mm.
 CG and PCG give iterations to a relative residual of $10^{-8}$.}
 \label{tab:results-objects}
 \begin{tabular}{lrrrrrr}
 \toprule
 Object & Free unknowns & $\bar u_x$ & $u_x^m$ & $\kappa_2(A)$ & CG & PCG\\
 \midrule
 Square & 1282 & 0.4936 & 0.4918 & $1.12\times10^{4}$ & 268 & 245\\
 Rectangle $2\times1$ & 2532 & 0.9936 & 0.9937 & $1.01\times10^{5}$ & 427 & 395\\
 Square with opening & 1246 & 0.5560 & 0.5883 & $1.30\times10^{4}$ & 283 & 256\\
 \bottomrule
 \end{tabular}
\end{table}

Doubling the length doubles the mean extension (a factor of $2.01$),
as expected for a plate in nearly uniform tension. The condition number
grows about nine times and CG needs 427 instead of 268 iterations.
The opening removes 4\% of the material but increases the mean extension
by 13\% and the midpoint extension by 20\%, because the load must flow
around it; the condition number rises by only 16\%.

Under refinement, the mean extension of the square and the rectangle
changes by less than 0.04\% at the last step, but that of the perforated
square changes by 0.66\%. The perforated square is therefore the most mesh
dependent, consistent with stresses concentrating at the corners of the
opening. All successive differences exceed the algebraic uncertainty
estimates.

The displacements scaled exactly with the load magnitude. Moving the same
total force to the lower or upper quarter of the right edge leaves the
mean extension almost unchanged ($0.4949$ against $0.4920$ mm for the
middle patch on the square), but bends the object: the right-edge midpoint
moves vertically by $\pm1.07$ mm on the square and $\pm4.42$ mm on the
rectangle. Clamping only the central half of the left edge increases the
mean extension by 22\% (square), 11\% (rectangle) and 25\% (perforated
square), and multiplies $\kappa_2(A)$ by 2.4 to 3.0.

The solver comparison repeats the square's pattern on every object.
CG, PCG, MINRES and GMRES reached the target on all 27 distinct systems.
Jacobi diverged on all of them, and damped Jacobi and steepest descent
exhausted the 20,000-iteration budget. Gauss--Seidel converged only on two
of the coarsest meshes: with the larger condition numbers of these
problems it becomes impractical.

\subsection{Validation and remaining limits}
Across the four refinement references, prescribed displacements are
exactly zero and true relative residuals are below $1.3\times10^{-13}$.
Normalized force and moment imbalances are below $4\times10^{-14}$,
counting reactions only at constrained components. Force imbalance is
normalized by the sum of applied nodal
force magnitudes; moment imbalance about the origin uses that force scale
times the square's diagonal length.  The separate roller-supported affine test gives a
relative nodal displacement error of $2.12\times10^{-15}$ against its
exact field; it does not provide an exact answer for the clamped benchmark.

The evidence therefore separates algebraic solution accuracy, mesh
dependence and verification against a known solution. It does not establish continuum accuracy or physical validity near the
clamp corners and the corners of the opening, where stresses concentrate.
No material or thickness comparison is made.
